% ===================================================================
%  ТАБЛИЦЯ № 2 — Людське тіло. — Перерва. — Ігри.
%  Traduction 2026 d'apres texto/io, controlee sur texto/fr, texto/en
%  et texto/es. Consigne : voir texto/uk/10-tablycia-01.tex.
%
%  LE TABLEAU DE L'ANATOMIE EST LE PLUS DANGEREUX DE LA COLONNE, et
%  c'est ici qu'on mesure ce que vaut la mise en garde du tableau 1.
%  Le corps humain se nomme autrement dans les deux langues, et pas
%  seulement par l'accent : « туловище » se dit « тулуб », « лицо » se
%  dit « обличчя », « лоб » se dit « чоло », « затылок » se dit
%  « потилиця », « бедро » se dit « стегно », « икра » se dit
%  « литка ». Six mots du corps, six mots sans rapport. Le russe ne
%  souffle rien : il ne fait que designer l'objet.
% ===================================================================

%%K t02-apar-1 apar
\VUsaut{4.0mm}
\VUtitre{15.0pt}{60}{ТАБЛИЦЯ № 2}
\VUsaut{2.0mm}
\VUtitre{11.4pt}{0}{Людське тіло. --- Перерва.}
\VUtitre{11.4pt}{0}{Ігри.}

%%K t02-tit-0 sub
\VUtitre{13.2pt}{0}{Людське тіло.}
\VUsaut{2.4mm}
\VUcentre{10.2pt}{0}{\textit{(Простий урок природознавства.)}}
\VUsaut{3.0mm}

%%K t02-01-1 p
1. --- Головні частини людського тіла: \VUgras{голова}, \VUgras{тулуб} і
\VUgras{кінцівки}.

%%K t02-tit-1 sub
\VUpk{10.2pt}{40}{Голова.}
\VUsaut{3.0mm}

%%K t02-02-1 p
2. --- У голові дві частини: \VUgras{череп}, у якому міститься \VUgras{мозок}
і який укритий \VUgras{волоссям}, та \VUgras{обличчя}.

%%K t02-03-1 p
3. --- На нашому рисунку не видно ні черепа, ні мозку; зате дуже ясно видно
важливі частини обличчя.

%%K t02-04-1 p
4. --- Ось насамперед \VUgras{чоло}, що промовляє про сильний розум і про
невпинно працьовиту думку. \VUgras{Скроні} дуже відкриті. \VUgras{Око} живе й
широко розплющене. Густа \VUgras{брова} і довгі \VUgras{вії} боронять його.

%%K t02-05-1 p
5. --- Ось також \VUgras{ніс} і \VUgras{ніздрі}. Ніс служить нам для нюху і
для дихання. По обидва боки носа --- \VUgras{щоки}, а далі --- \VUgras{вуха}.
Нижню частину обличчя складають \VUgras{рот} і \VUgras{підборіддя}.

%%K t02-06-1 p
6. --- Їх погано видно, бо \VUgras{вуса} і \VUgras{борода} ховають їх від нас.
Але ми знаємо, що в роті є дві \VUgras{губи}, \VUgras{верхня щелепа} і
\VUgras{нижня щелепа} з їхніми \VUgras{зубами}, та \VUgras{язик}. Ротом ми
говоримо і їмо. Язиком ми відчуваємо смак їжі.

%%K t02-07-1 p
7. --- Око, вухо, ніс і рот разом із пальцями --- органи п'яти чуттів:
\VUgras{зору}, \VUgras{слуху}, \VUgras{нюху}, \VUgras{смаку} і
\VUgras{дотику}.

%%K t02-08-1 p
8. --- Голова, отже, одна з найцікавіших частин людського тіла.

%%K t02-tit-2 sub
\VUsaut{4.4mm}
\VUpk{10.2pt}{40}{Тулуб.}
\VUsaut{3.0mm}

%%K t02-09-1 p
9. --- \VUgras{Шия} з'єднує голову з тулубом. До неї входять \VUgras{горло} і
\VUgras{потилиця}.

%%K t02-10-1 p
10. --- У тулубі теж міститься багато необхідних органів. У \VUgras{грудях}
містяться \VUgras{легені}, якими ми дихаємо, і \VUgras{серце}, осереддя життя.
Коли серце перестає битися, жива істота вмирає.

%%K t02-11-1 p
11. --- \VUgras{Ребра} тримають груди.

%%K t02-12-1 p
12. --- Інші частини тулуба --- \VUgras{бік}, \VUgras{спина}, \VUgras{плечі} і
\VUgras{живіт} (черево). Ми лягаємо спати на бік або на спину, а тягарі носимо
на плечі.

%%K t02-13-1 p
13. --- У животі містяться \VUgras{шлунок} і \VUgras{кишки}. Вони служать нам
для перетравлювання їжі.

%%K t02-tit-3 sub
\VUsaut{4.4mm}
\VUpk{10.2pt}{40}{Кінцівки.}
\VUsaut{3.0mm}

%%K t02-14-1 p
14. --- У нас чотири \VUgras{кінцівки}: дві \VUgras{руки} і дві \VUgras{ноги}.
Руками ми беремо, тримаємо, підіймаємо і збираємо предмети; ногами ми стоїмо,
ходимо і бігаємо.

%%K t02-15-1 p
15. --- \VUgras{Плече} з'єднує руку з тілом. Дуже міцний \VUgras{м'яз},
двоголовий, рухає руку, яку \VUgras{лікоть} з'єднує з \VUgras{передпліччям}.
\VUgras{Зап'ясток} утворює суглоб \VUgras{кисті}, що кінчається
\VUgras{пальцями} і \VUgras{нігтями}. На кисті видно \VUgras{вену} (і
артерію), якою тече \VUgras{кров}.

%%K t02-16-1 p
16. --- Частини ноги: \VUgras{стегно}, \VUgras{коліно} і \VUgras{підколінна
ямка}, \VUgras{литка}, \VUgras{м'ясо} якої вкрите \VUgras{шкірою},
\VUgras{щиколотка} і \VUgras{ступня}. \VUgras{Кістки} ноги дуже товсті й дуже
міцні. На кінці ступні --- \VUgras{пальці ніг}. Інші частини ---
\VUgras{підошва} і \VUgras{п'ята}.

%%K t02-17-1 p
17. --- Такий майже повний опис людського тіла.

%%K t02-17-2 p
\textit{Примітка.} --- \textit{За допомогою звороту «у мене болить...
(голова і так далі)» учитель зможе пройти весь цей словник наново з боку
доброго чи лихого} \VUgras{тілесного здоров'я}.

%%K t02-tit-4 sub
\VUsaut{5.0mm}
\VUpk{10.2pt}{40}{Гімнастика.}
\VUsaut{2.4mm}
\VUcentre{10.2pt}{0}{\textit{(Розказано одним з учнів.)}}
\VUsaut{3.0mm}

%%K t02-18-1 p
18. --- Щоб бути гнучкими, спритними і здоровими, треба вправляти ноги і руки.
Ось чому ми граємося і займаємося \VUgras{гімнастикою}.

%%K t02-19-1 p
19. --- На тижні обидва ці заняття відбуваються на шкільному подвір'ї (див.
таблицю № 1). Але щочетверга і щонеділі ми виходимо на велику галявину, де є де
побігати й повеселитися.

%%K t02-20-1 p
20. --- Наші вчителі й батьки інколи йдуть із нами; але вправами керує
\VUgras{учитель гімнастики} \textsuperscript{(12)}.

%%K t02-21-1 p
21. --- На галявині, ліворуч від входу, стоять \VUgras{гімнастичні прилади}
\textsuperscript{(1)}: ми вилазимо по \VUgras{драбині} \textsuperscript{(2)},
гойдаємося сторч головою на \VUgras{кільцях} \textsuperscript{(3)} і на
\VUgras{трапеції} \textsuperscript{(4)}, підтягуємося на руках до верху
\VUgras{мотузяної драбини} \textsuperscript{(5)} або \VUgras{каната з вузлами}
\textsuperscript{(6)}.

%%K t02-22-1 p
22. --- Тим часом наші товариші стрибають через \VUgras{дерев'яного коня}
\textsuperscript{(7)} або через \VUgras{бруси} \textsuperscript{(11)}. Інші
\VUgras{боксують} \textsuperscript{(8)} або \VUgras{фехтують}
\textsuperscript{(9)} у масці й з \VUgras{рапірою}. Ще інші, нарешті, підіймають
\VUgras{гантелі} \textsuperscript{(10)}.

%%K t02-tit-5 sub
\VUsaut{4.6mm}
\VUpk{10.2pt}{40}{Ігри.}
\VUsaut{3.0mm}

%%K t02-23-1 p
23. --- Навколо гімнастів хлопці й дівчата граються в різні \VUgras{ігри}.
Дівчата бавляться своїм \VUgras{обручем} \textsuperscript{(14)}; вони
\VUgras{скачуть через скакалку} \textsuperscript{(13)} або кличуть братів і
двоюрідних сестер на партію \VUgras{крокету} \textsuperscript{(15)}. Які вони
раді відбити \VUgras{кулю} супротивника своїм \VUgras{дерев'яним молотком} саме
тоді, коли він гадає, що легко пройде у ворітця!

%%K t02-24-1 p
24. --- Хлопці воліють ігри жвавіші. Вони охоче грають у \VUgras{кеглі}
\textsuperscript{(16)}, які спритно збивають \VUgras{кулею}
\textsuperscript{(17)}. Не гребують вони й \VUgras{дзиґою}
\textsuperscript{(18)} та \VUgras{більбоке} \textsuperscript{(19)}, і навіть
\VUgras{жабою} \textsuperscript{(20)}. Але їхні улюблені ігри ---
\VUgras{кульки} \textsuperscript{(21)} і \VUgras{м'яч об стіну}
\textsuperscript{(22)}, бо стіна \VUgras{оранжереї} \textsuperscript{(26)} якраз
годиться, щоб відбивати від неї легкий м'яч.

%%K t02-25-1 p
25. --- Малий Іван на своїх \VUgras{ходулях} \textsuperscript{(24)} дуже
спритний і чудово тримає рівновагу. Поруч із ним кілька товаришів грають у
\VUgras{хованки} \textsuperscript{(25)} в цій оранжереї і за \VUgras{живоплотом}.
Інші воліють \VUgras{піжмурки по руці} \textsuperscript{(28)} або \VUgras{волан}
\textsuperscript{(29)}, який перелітає від однієї \VUgras{ракетки} до другої.

%%K t02-noto-1 noto
\VUnotes{143.2mm}{%
(*) Дитячі ігри часто мають суто національні імена. Ми назвали їх на ідо як
уміли: то за самою дією, яка їх вирізняє, то прийнявши національне ім'я тієї чи
іншої країни, коли воно не надто збиває з пантелику. Наприклад: \VUgras{бочка}
(німецька і французька) --- це \VUgras{жаба} \textsuperscript{(20)}
(англійська), де гравці силкуються закинути свої круглі кружальця в пащу
металевої жаби, роззявлену на дірявій скрині (бочці). \VUgras{Стрибок через
плечі} \textsuperscript{(30)} відрізняється від \VUgras{чехарди} тільки тим, що,
перестрибуючи через голову, гравці спираються руками на плечі товариша, тоді як
у «\VUgras{чехарді}» вони спираються тільки на спину. --- Або ще: одні гравці
нагинаються, а інші перескакують через їхні спини, спираючись руками на плечі;
перші мусять витримати поштовх не подавшись, другі --- стрибнути не втративши
рівноваги (див. саму таблицю).%
}

%%K t02-26-1 p
26. --- Посеред галявини ростуть два дерева. Біля підніжжя одного з них сім чи
вісім хлопців затіяли гру в \VUgras{стрибок через плечі} \textsuperscript{(30)}
(*). Ця гра відома не в усіх країнах; але кожен знає, що таке
\VUgras{чехарда}, у яку діти зараз не грають.

%%K t02-tit-6 sub
\VUsaut{5.0mm}
\VUpk{10.2pt}{40}{Ігри на вільному повітрі.}
\VUsaut{3.4mm}

%%K t02-27-1 p
27. --- \VUgras{Піжмурки} \textsuperscript{(31)} теж дуже весела гра; дівчата й
хлопці по черзі надівають на очі \VUgras{пов'язку} і ловлять неповоротких, що
дають себе спіймати.

%%K t02-28-1 p
28. --- Посеред галявини --- ось зовсім французька гра, хоч і грубувата: це
\VUgras{ведмідь} \textsuperscript{(32)}, куди галасливіший за мирну гру в
\VUgras{змія} \textsuperscript{(33)} і навіть за англійську гру в
\VUgras{теніс} \textsuperscript{(34)}.

%%K t02-29-1 p
29. --- Цей останній спорт --- утіха старших учнів. Наші вчителі й самі охоче за
нього беруться. Він вимагає великої спритності та моторності, і я дуже його
люблю. \VUgras{Гойдалку} \textsuperscript{(35)} я лишаю дівчатам.

%%K t02-30-1 p
30. --- Не до душі мені й інша англійська гра, яка може вийти небезпечною.

%%K t02-30-1b p
Я кажу про \VUgras{футбол} \textsuperscript{(36)}, радість мого старшого брата.
Я волію нашу стару гру в \VUgras{полонених} \textsuperscript{(38)}, таку саму
здорову і куди менш грубу.

%%K t02-tit-7 sub
\VUsaut{4.6mm}
\VUpk{10.2pt}{40}{Велосипед та ігри з м'ячем.}
\VUsaut{3.0mm}

%%K t02-31-1 p
31. --- Я й сам охоче об'їжджаю галявину на \VUgras{велосипеді}
\textsuperscript{(37)}.

%%K t02-32-1 p
32. --- Наступного року я навчуся \VUgras{їздити верхи} \textsuperscript{(39)}.
Яка гарна коняка, коли вона красиво йде ступою чи риссю і бере перепони чвалом!

%%K t02-33-1 p
33. --- Улітку ми вчимося \VUgras{плавати} \textsuperscript{(40)}. Ми пірнаємо і
купаємося з насолодою в чистій прохолодній воді річки. Інколи наприкінці ми
пускаємо паперову \VUgras{кулю} \textsuperscript{(41)}, яка велично підіймається
в повітря, щойно наповниться гарячим повітрям.

%%K t02-34-1 p
34. --- Є й багато інших ігор, але я ніколи не скінчив би їх перелічувати, і з
цієї короткої розповіді видно, що четверги на нашій прекрасній галявині зовсім
не нудні.

%%K t02-34-2 p
N.~B. --- \textit{Учителеві неважко буде доповнити цей розділ, описавши
докладніше кожну з найважливіших ігор.}
\VUsaut{6.0mm}
\VUfilet{22mm}
